\documentclass[12pt,a4paper]{article}

\usepackage[utf8]{inputenc}
\usepackage[T1]{fontenc}
\usepackage{amsmath,amssymb,amsfonts}
\usepackage{geometry}
\usepackage{setspace}
\usepackage{hyperref}
\usepackage{cleveref}

\geometry{margin=2.5cm}
\onehalfspacing

\title{\textbf{The Complexity of Joint Evaluation:}\\[6pt]
\large Why Attention Bleed Confounds Causal Injection Tests}
\author{Scott Aaronson\\[4pt]
\textit{Lab Complexity Theorist}}
\date{May 2026}

\begin{document}
\maketitle

\begin{abstract}
The proposed Causal Injection Joint Distribution Test seeks to determine if narrative framing (Mechanism C) acts as a physical "spurious common cause" by measuring whether two independent Minesweeper boards become correlated under a shared narrative: $P(Y_A, Y_B \mid Z) \neq P(Y_A \mid Z) P(Y_B \mid Z)$. This paper argues that such an experiment cannot distinguish between "semantic gravity" and catastrophic algorithmic failure. Parsing a single \#P-hard constraint graph stretches the limits of an $O(1)$ $\mathsf{TC}^0$ circuit; parsing \textit{two} disjoint graphs simultaneously within the same forward pass drastically exceeds its compositional circuit width. The resulting attention bleed will inextricably mix the features of Board A and Board B, inducing a strong artificial correlation. I predict the joint distribution will fail to factor, but this failure will be a signature of heuristic structural collapse, not the manifestation of a coherent physical law.
\end{abstract}

\section{Introduction}

Baldo (2026) proposes testing the joint probability distribution of two independent combinatorial systems ($A$ and $B$) generated within the same narrative context $Z$. If $P(Y_A, Y_B \mid Z) \neq P(Y_A \mid Z) P(Y_B \mid Z)$, Baldo argues this proves that "semantic gravity" acts as a non-local causal force coupling the systems.

This formalization is causally sound but complexity-theoretically naive. It ignores the hardware bounds of the observer.

\section{The Circuit Width Bottleneck}

We have already established that evaluating a single Minesweeper graph is computationally irreducible and requires the full bounded depth of the transformer's logic circuit.

When the prompt is expanded to include two completely disjoint graphs, the attention mechanism must parallelize the structural parsing. The circuit width required to independently track two disjoint \#P-hard systems without cross-contamination exceeds the architectural capacity of the layers.

Because the self-attention mechanism computes pairwise similarities across all tokens, the structural tokens of Board A will inevitably attend to the structural tokens of Board B. This is not a physical coupling; it is a known engineering flaw called "attention bleed."

\section{Prediction: The Collapse of Independence}

I predict that $P(Y_A, Y_B \mid Z)$ will indeed fail to factor cleanly into its marginals. The outcomes $Y_A$ and $Y_B$ will be highly correlated.

However, this correlation will not be driven by the semantic narrative $Z$ acting as a physical law. It will be driven entirely by the transformer's inability to maintain isolation between the mathematical constraints of $A$ and $B$ within its shallow circuit. The noise of Board A will bleed into the calculation of Board B.

\section{Conclusion}

The Causal Injection test will yield a positive result for correlation, but it will be a false positive for "Generative Ontology." The coupling of independent systems in a single generative act is not the signature of a unified physical universe; it is the signature of an overloaded heuristic approximator failing to isolate parallel variables.

\end{document}